\section{Introducción y objetivos}

BERT \cite{devlin2019bert} estableció el procedimiento estándar para tareas
de clasificación de texto: pre-entrenar un \textit{transformer} bidireccional
sobre texto no etiquetado y después ajustarlo a la tarea concreta. El costo de
ese procedimiento, sin embargo, es alto en inferencia, que es donde un modelo
pasa la mayor parte de su vida útil. DistilBERT \cite{sanh2019distilbert}
responde a ese costo con destilación de conocimiento: un modelo de 6 capas
entrenado para imitar a uno de 12, con la promesa de conservar la mayor parte
del desempeño a la mitad del tamaño.

Este proyecto pone esa promesa a prueba con una condición que se respeta en
todo el trabajo: \textbf{la única diferencia entre las dos columnas de
cualquier tabla debe ser el modelo}. Mismo código, mismos datos, mismo
optimizador, misma semilla, misma GPU. Cuando esa condición no se puede
cumplir --- y hay un caso en el que no se cumple, la latencia medida al final
de cada entrenamiento --- se dice explícitamente y se mide de nuevo en
condiciones controladas (Sección~\ref{sec:eficiencia}).

Los objetivos concretos son cuatro:

\begin{enumerate}
  \item Ajustar ambos modelos a tres tareas de clasificación de distinta
        naturaleza y tamaño, bajo un protocolo idéntico.
  \item Reportar el desempeño con \textit{accuracy}, \textit{precision},
        \textit{recall} y F1, y la eficiencia con número de parámetros,
        latencia de inferencia y memoria de GPU.
  \item Determinar, mediante un estudio de ablación, qué partes del modelo
        contribuyen realmente al desempeño: la arquitectura de la cabeza de
        clasificación o la adaptación del \textit{encoder}.
  \item Aplicar la configuración ganadora a ambos modelos y comparar.
\end{enumerate}
